\section{Auxiliary Supervision}

The autoencoder architectures described previously are trained primarily
to preserve information required for reconstruction of the dynamic BOLD
signal. Reconstruction alone, however, does not explicitly require the
latent representation to retain information associated with disease state
or pathological biomarkers. Auxiliary supervision was therefore introduced
to investigate whether clinically relevant information could be encouraged
within the latent space while maintaining the ability of the autoencoder
to reconstruct dynamic brain activity.

Two forms of auxiliary supervision were investigated independently:
diagnostic classification and biomarker regression. Diagnostic supervision
used the HC, MCI, and AD class labels, while regression supervision used
continuous amyloid-$\beta$ and tau measurements. In both cases, auxiliary
heads were attached directly to the latent representation, as illustrated
in Figure~\ref{fig:auxiliary_heads}. Gradients from the auxiliary objectives
therefore propagated through the prediction head and encoder, encouraging
the latent representation to support the corresponding supervised task,
while the reconstruction objective continued to constrain preservation of
the original BOLD signal.

\thesisfigure{architectures/auxiliaryHeads}{Conceptual addition of auxiliary supervision to the autoencoder framework}{fig:auxiliary_heads}{\textwidth}





\subsection{Diagnostic Classification Head}

Diagnostic classification was investigated as a source of categorical
disease-related supervision for the latent representation. The objective
of the classification head was to encourage the encoder to retain features
associated with diagnostic status while preserving the information required
for reconstruction of the dynamic BOLD signal. The classification task was
therefore treated as an auxiliary objective rather than as the primary
purpose of the model.

The time-preserving latent representation,
\[
\mathbf{Z}\in\mathbb{R}^{T\times L},
\]
was flattened into a subject-level feature vector,
\[
\mathbf{z}_{\mathrm{flat}}\in\mathbb{R}^{TL},
\]
before being provided to the classification head. This allows information
from all latent dimensions and time points to contribute jointly to the
subject-level diagnostic prediction. 

Two classification-head architectures were considered: a linear classifier
and a small multilayer perceptron (MLP). The linear classifier consisted of
a single affine transformation,
\begin{equation}
    \mathbf{o}
    =
    \mathbf{W}_{\mathrm{cls}}
    \mathbf{z}_{\mathrm{flat}}
    +
    \mathbf{b}_{\mathrm{cls}},
\end{equation}
where
\[
\mathbf{o}\in\mathbb{R}^{C}
\]
denotes the classifier output and $C$ is the number of diagnostic classes.
For the HC, MCI, and AD classification task, $C=3$. This head provides the
lowest-capacity classification variant and tests whether diagnostic
information can be extracted directly from the flattened latent
representation through a linear mapping.

The MLP classifier introduces an additional nonlinear hidden transformation,
\begin{equation}
    \mathbf{h}_{\mathrm{cls}}
    =
    \operatorname{ReLU}
    \left(
        \mathbf{W}_{1}
        \mathbf{z}_{\mathrm{flat}}
        +
        \mathbf{b}_{1}
    \right),
\end{equation}
followed by a second linear transformation,
\begin{equation}
    \mathbf{o}
    =
    \mathbf{W}_{2}
    \mathbf{h}_{\mathrm{cls}}
    +
    \mathbf{b}_{2}.
\end{equation}
The dimensional progression of this head can therefore be represented as
\[
TL
\rightarrow
H_{\mathrm{cls}}
\rightarrow
C,
\]
where $H_{\mathrm{cls}}$ denotes the dimensionality of the hidden layer.
The nonlinear transformation increases the representational capacity of
the auxiliary classifier and allows interactions between latent features
to contribute to diagnostic prediction. The implemented classification
heads therefore differ primarily in whether diagnostic information is
modelled using a direct linear mapping or a small nonlinear network.


The classification objective was computed using cross-entropy loss, as
defined in the Training Objectives section. The classification loss was
combined with the reconstruction objective through a weighting coefficient
$\lambda_{\mathrm{cls}}$,
\begin{equation}
    \mathcal{L}_{\mathrm{total}}
    =
    \mathcal{L}_{\mathrm{reconstruction}}
    +
    \lambda_{\mathrm{cls}}
    \mathcal{L}_{\mathrm{cls}}.
\end{equation}

The classification head was trained jointly with the autoencoder. As a
result, gradients from the classification objective propagated through the
classification head and into the encoder, allowing diagnostic supervision
to influence the latent representation. At the same time, the
reconstruction objective continued to constrain the encoder-decoder model
to preserve the dynamic BOLD activity. The coefficient
$\lambda_{\mathrm{cls}}$ therefore controls the relative contribution of
diagnostic supervision to the overall training objective.

The classification experiments were designed to evaluate whether adding
diagnostic supervision could encourage the latent representation to encode
more class-related information without degrading reconstruction quality.
The auxiliary models were therefore evaluated using the same reconstruction
and dynamic-connectivity metrics as the preceding autoencoder architecture
experiments, while classification performance was used as an additional
measure of how effectively diagnostic information could be extracted from
the learned latent representation.




\subsection{Biomarker Regression Heads}

Biomarker regression was investigated as a source of continuous
disease-related supervision for the latent representation. In contrast to
the diagnostic classification objective, which uses categorical disease
labels, the regression objective provides supervision based on continuous
amyloid-$\beta$ and tau measurements. The purpose of these heads was to
encourage the encoder to retain information associated with Alzheimer's
disease pathology while preserving the information required for
reconstruction of the dynamic BOLD signal.

The regression heads operate directly on the time-preserving latent
representation,
\[
    \mathbf{Z}\in\mathbb{R}^{T\times L}.
\]
Rather than flattening the complete latent representation, as performed for
the diagnostic classification head, the regression architectures aggregate
or model the temporal dimension before producing a subject-level biomarker
prediction. This preserves the temporal organisation of the latent
representation within the auxiliary head and allows different assumptions
about the distribution of biomarker-related information over time to be
investigated.

Three regression-head architectures were considered: an average-pooling
head, a temporal-attention pooling head, and a temporal-convolutional head.
These architectures progressively introduce learned temporal processing,
ranging from uniform aggregation across time to explicit modelling of local
temporal patterns.

\subsubsection{Average-Pooling Regression Head}

The average-pooling head provides the simplest regression architecture. The
latent representation is first averaged over the temporal dimension,
producing a subject-level latent vector
\[
    \bar{\mathbf{z}}
    =
    \frac{1}{T}
    \sum_{t=1}^{T}
    \mathbf{z}_{t},
    \qquad
    \bar{\mathbf{z}}\in\mathbb{R}^{L}.
\]
A linear transformation then maps this representation to the predicted
biomarker value,
\begin{equation}
    \hat{\mathbf{y}}
    =
    \mathbf{W}_{\mathrm{reg}}
    \bar{\mathbf{z}}
    +
    \mathbf{b}_{\mathrm{reg}}.
\end{equation}

This architecture assigns equal importance to every time point and therefore
does not explicitly model temporal variation within the latent
representation. It provides a low-capacity baseline for determining whether
biomarker-related information is represented as an overall subject-level
pattern across the latent trajectories.

\thesisfigure
    {architectures/avgRegHead}
    {Average-pooling biomarker regression head. The temporal dimension of
    the latent representation is averaged before a linear projection to the
    biomarker prediction.}
    {fig:avg_regression_head}
    {0.75\textwidth}


\subsubsection{Temporal-Attention Pooling Regression Head}

Uniform temporal averaging assumes that all time points contribute equally
to biomarker prediction. To allow the model to learn whether particular
parts of the latent trajectory are more informative, a temporal-attention
pooling head was additionally considered.

For each time point, the corresponding latent vector
$\mathbf{z}_{t}\in\mathbb{R}^{L}$ is mapped to a scalar attention score,
\begin{equation}
    e_t
    =
    \mathbf{w}_{a}^{\mathsf{T}}
    \mathbf{z}_{t}
    +
    b_a.
\end{equation}
The scores are normalised across time,
\begin{equation}
    \alpha_t
    =
    \frac{\exp(e_t)}
    {\sum_{j=1}^{T}\exp(e_j)},
\end{equation}
such that
\[
    \sum_{t=1}^{T}\alpha_t=1.
\]
A subject-level representation is then constructed as the weighted sum of
the latent vectors,
\begin{equation}
    \mathbf{z}_{\mathrm{att}}
    =
    \sum_{t=1}^{T}
    \alpha_t\mathbf{z}_{t},
\end{equation}
where
$\mathbf{z}_{\mathrm{att}}\in\mathbb{R}^{L}$.
The pooled representation is subsequently mapped to the biomarker
prediction using a linear regression layer,
\begin{equation}
    \hat{\mathbf{y}}
    =
    \mathbf{W}_{\mathrm{reg}}
    \mathbf{z}_{\mathrm{att}}
    +
    \mathbf{b}_{\mathrm{reg}}.
\end{equation}

Unlike average pooling, this architecture learns the relative contribution
of individual time points to the auxiliary prediction. It therefore tests
whether biomarker-related information is more strongly represented during
particular portions of the latent trajectory rather than being distributed
uniformly across the scan.

\thesisfigure
    {architectures/tempAttPoolRegHead}
    {Temporal-attention pooling biomarker regression head. Learned attention
    weights are used to construct a weighted temporal summary of the latent
    representation before biomarker prediction.}
    {fig:temporal_pool_regression_head}
    {0.85\textwidth}

\subsubsection{Temporal-Convolutional Regression Head}

The temporal-convolutional head was introduced to investigate whether local
patterns in the latent trajectories contain biomarker-related information
that cannot be captured through direct temporal aggregation. In contrast to
the preceding pooling approaches, temporal convolutions transform the
latent sequence before it is reduced to a subject-level representation.

Internally, the latent representation is arranged as
\[
    \mathbf{Z}\in\mathbb{R}^{L\times T},
\]
such that the $L$ latent dimensions form the input channels and convolution
is performed along the temporal dimension. The regression head applies a
one-dimensional convolution with kernel size three, followed by a GELU
activation. When a hidden representation is used, a second temporal
convolution and GELU activation are applied,
\[
    (L,T)
    \rightarrow
    (H_{\mathrm{reg}},T)
    \rightarrow
    (H_{\mathrm{reg}},T),
\]
where $H_{\mathrm{reg}}$ denotes the number of hidden feature channels.
Padding is used such that the temporal length is retained by the
convolutional transformations.

Adaptive average pooling subsequently reduces the temporal dimension,
\[
    (H_{\mathrm{reg}},T)
    \rightarrow
    (H_{\mathrm{reg}},1),
\]
and the resulting feature vector is mapped to the biomarker prediction by a
linear output layer. The complete transformation can therefore be
summarised as
\[
    (L,T)
    \rightarrow
    (H_{\mathrm{reg}},T)
    \rightarrow
    (H_{\mathrm{reg}},T)
    \rightarrow
    H_{\mathrm{reg}}
    \rightarrow
    \hat{\mathbf{y}}.
\]

Whereas average pooling tests for temporally distributed biomarker
information and attention pooling allows individual time points to receive
different importance, the convolutional head explicitly models local
relationships between neighbouring time points. It therefore provides a
complementary test of whether short-range temporal patterns within the
latent representation contribute to biomarker prediction.

\thesisfigure
    {architectures/convRegHead}
    {Temporal-convolutional biomarker regression head. Temporal
    convolutions extract local patterns from the latent trajectories before
    adaptive average pooling and linear biomarker prediction.}
    {fig:conv_regression_head}
    {0.9\textwidth}

For each architecture, the regression output was trained against the
corresponding amyloid-$\beta$ or tau measurement using the Smooth L1
objective described in the Training Objectives section. When multiple
biomarker prediction objectives were used, their contributions were
combined through the regression loss before being weighted by the
auxiliary-loss coefficient $\lambda_{\mathrm{pred}}$.

The three regression architectures were selected to provide distinct levels
of temporal modelling while retaining relatively compact auxiliary heads.
The average-pooling model contains no learned temporal aggregation, the
temporal-attention model learns the importance of individual time points,
and the temporal-convolutional model learns local temporal features before
aggregation. Their comparison therefore evaluates not only whether
biomarker supervision can influence the latent representation, but also
whether the manner in which temporal latent information is exposed to the
auxiliary task affects preservation of dynamic BOLD reconstruction.
